\section{Related Work}

\textbf{AutoResearch systems and evaluation targets.}
AutoResearch systems increasingly automate hypothesis formation, experimentation, and artifact production, from end-to-end scientific workflows and agentic tree search to focused, compute-bounded model-training loops~\citep{lu2024aiscientist,yamada2025aiscientistv2,karpathy2026autoresearch}.
Recent systems further explore persistent project state, hypothesis-tree refinement, cross-task skill accumulation, and continuously evolving multi-agent workflows~\citep{chen2026toward,jin2026hypothesistree,kim2026haste,zhu2026evomaster}.
Benchmarks isolate complementary slices of this research loop: RE-Bench and PostTrainBench target open-ended AI R\&D and autonomous model post-training~\citep{wijk2025rebench,rank2026posttrainbench}; MLGym, AIRS-Bench, MLRC-Bench, and ResearchGym study research-oriented problem solving and end-to-end ML projects~\citep{nathani2025mlgym,lupidi2026airs,zhang2025mlrcbench,garikaparthi2026researchgym}; and PaperBench and NatureBench emphasize reproducing papers or published scientific results~\citep{starace2025paperbench,wang2026naturebench}.
Within executable MLE, MLE-Bench evaluates agents on end-to-end Kaggle-style competitions~\citep{chan2024mlebench}, whereas MLS-Bench targets improving ML components in ways that generalize across controlled settings and scales~\citep{lyu2026mlsbench}.
We evaluate OpenMLE on MLE-Bench Lite as its primary MLE benchmark and on NatureBench Lite as a focused test of transfer to broader scientific AutoResearch.

\textbf{Executable MLE environments and scalable task resources.}
Classical AutoML optimizes over predefined model and pipeline spaces~\citep{thornton2013autoweka,feurer2015autosklearn,olson2016tpot,erickson2020autogluon}, whereas LLM-based MLE agents operate through open-ended, code-mediated experimentation.
Building on this benchmark lineage, MLAgentBench and DSBench likewise require agents to write, execute, debug, and submit solutions on realistic ML and data-science tasks~\citep{huang2023mlagentbench,jing2024dsbench}.
Complementary efforts expand the underlying execution and training infrastructure: MLE-Dojo standardizes interactive MLE environments~\citep{qiang2025mledojo}, while MLE-Smith and SandMLE scale the construction of executable training tasks~\citep{qiang2025mlesmith,zhou2026sandmle}.
These resources make execution feedback available at increasing scale, but executable tasks alone do not specify how experience should be transformed into reusable model capabilities or composed at inference time.
\openmlegym{} builds on this line by unifying scalable task environments with isolated execution, structured feedback, and task-specific evaluators as a shared substrate for post-training, search, and evaluation.

\textbf{Inference-time scaffolds and evolutionary search.}
Given executable evaluators, inference-time scaffolds amplify frozen models by allocating trials, maintaining search state, and selecting or transforming candidate programs.
MLE systems instantiate this idea through multi-agent decomposition~\citep{trirat2024automlagent,gandhi2024budgetmlagent,li2025autokaggle,fang2025mlzero}, structured search~\citep{jiang2025aide,chi2024sela}, and targeted refinement or domain knowledge~\citep{liu2025ml,ou2025automind,nam2025mlestar,du2025automlgen,zhang2025deepanalyze}.
AIRA-style analyses further identify search policy, operator quality, throughput, and ideation diversity as key performance factors~\citep{toledo2025airesearchagents,audranreiss2025ideationdiversity,hambardzumyan2026aira_2}.
More generally, program-evolution systems use language models as mutation operators and executable evaluators as selection signals~\citep{novikov2025alphaevolve,lange2025shinkaevolve,cemri2026adaevolve,ye2026evaluation}, while ThetaEvolve and TTT-Discover update problem-solving behavior from test-time feedback~\citep{wang2025thetaevolve,yuksekgonul2026learning}.
These methods establish the value of structured search, but typically rely on transformation behavior supplied by the frozen backbone or its prompts; \openmleevo{} instead composes operators explicitly trained for the same roles used during search.

\textbf{Learning MLE agents from executable experience.}
A complementary line internalizes verifiable outcomes through supervised or reinforcement-learning updates rather than retaining all improvement logic in an external scaffold.
RLVR has produced strong learned reasoning capabilities in mathematics and coding~\citep{guo2025deepseek,team2025kimi,zeng2025rlve} and has increasingly been extended to long-horizon agent tasks~\citep{team2026kimi,zeng2026glm,minimax2026m27}.
MLE-specific efforts similarly learn from executable task rewards~\citep{liu2025mlagent,li2025mlerl,yang2025rlmleagents,cai2026acegrpo}, establishing that MLE experience can be internalized by the model.
However, these efforts expose different subsets of the task, environment, training, evaluation, and model artifacts needed for a reproducible end-to-end workflow, and their learned policies are not uniformly organized around an interface shared with long-horizon search (Appendix Table~\ref{tab:mle-open-release-matrix}).
\openmleerl{} follows this learning direction but trains reusable \textsc{Draft}, \textsc{Improve}, \textsc{Debug}, and \textsc{Crossover} transformations whose interfaces match the operations invoked by \openmleevo{}.

\textbf{AI-for-AI and trainable improvers.}
AI-for-AI broadens executable improvement from optimizing task solutions to improving the models, operators, and harnesses that generate future solutions~\citep{jiang2026selfimprovingagents}.
Search--learning systems begin to return execution experience to the generator, as in test-time policy updates and schemes that alternate evolutionary search with hindsight fine-tuning~\citep{wang2025thetaevolve,yuksekgonul2026learning,pourcel2025soar}.
Related work extends improvement from candidate programs to agent and harness design~\citep{hu2024adas,zhang2026hyperagents,lee2026metaharness,zhang2026aevo}.
OpenMLE contributes at the interface of these directions: verified evolutionary experience post-trains the same program-transformation operators that subsequently govern evolutionary search, creating a meta-evolutionary coupling between learning and inference.
Relative to a public landscape that often exposes individual or partial components of the workflow, OpenMLE is organized as an open stack spanning task packages, sandbox execution, operator training, search, evaluation, and released model weights.
This coupling provides a concrete testbed for studying progress toward RSI in executable MLE, rather than a claim that OpenMLE realizes general, autonomous recursive self-improvement.
